\chapter{Vitamin E}

\section*{What Is This Test?}
Vitamin E testing measures serum or plasma alpha-tocopherol, the principal bioactive form of vitamin E in humans. Vitamin E is a fat-soluble antioxidant that protects polyunsaturated fatty acids in cell membranes from lipid peroxidation. The test is indicated for suspected deficiency (fat malabsorption, abetalipoproteinemia, cystic fibrosis, cholestatic liver disease, premature infants) or toxicity (high-dose supplementation). Routine population screening is not recommended. Vitamin E deficiency causes neurologic deficits (ataxia, peripheral neuropathy, retinopathy) and hemolytic anemia in neonates.

\section*{Alternative Names}
Alpha-Tocopherol; Vitamin E (Tocopherol); Serum Vitamin E; Plasma Tocopherol

\section*{Principle}
High-performance liquid chromatography (HPLC) with UV or fluorescence detection or LC-MS/MS quantifies alpha-tocopherol from serum. Lipid-standardized alpha-tocopherol (tocopherol/total lipid ratio) corrects for hyperlipidemia. Chromatographic methods also resolve beta-, gamma-, and delta-tocopherol isoforms if needed.

\section*{History}
Vitamin E was discovered by Evans and Bishop in 1922 as a dietary factor essential for reproduction in rats. Tocopherol structure was elucidated in the 1930s. Clinical deficiency in humans (especially from fat malabsorption disorders) was characterized in the 1960s--1980s. Controversial antioxidant-heart disease trials in the 1990s and 2000s found no benefit of vitamin E supplementation in cardiovascular disease.

\section*{Why This Test Is Ordered}
\begin{itemize}
  \item Suspected vitamin E deficiency in fat malabsorption syndromes: cystic fibrosis, abetalipoproteinemia, cholestatic liver disease, short bowel syndrome, Crohn's disease
  \item Evaluation of premature infants with hemolytic anemia or retinopathy
  \item Neurologic symptoms (spinocerebellar ataxia, peripheral neuropathy, myopathy) seen in chronic malabsorption
  \item Monitoring patients on long-term parenteral nutrition without adequate vitamin E
  \item Suspected vitamin E toxicity from megadose supplementation
\end{itemize}

\section*{Primary vs Secondary Test}
Secondary test ordered when clinical suspicion arises, especially in fat malabsorption syndromes.

\section*{How This Test Is Performed}
Venous blood drawn in serum separator tube (gold) or EDTA (lavender), preferably fasting. Specimen protected from light and hemolysis. Results available within 3--7 days. Reporting often includes alpha-tocopherol and lipid-standardized ratio.

\section*{What Preparation Is Needed}
Overnight fast preferred. Avoid vitamin E supplements for 24 hours before testing if measuring steady state on dietary intake. Stop multivitamins before testing.

\section*{Contraindications and Precautions}
Interpret with serum lipids because alpha-tocopherol circulates bound to lipoproteins. Low total cholesterol (familial hypobetalipoproteinemia, severe malabsorption) leads to low alpha-tocopherol despite adequacy; lipid-standardized ratio corrects. Hemolysis and light exposure degrade specimen. Heparin and anticonvulsants may lower vitamin E.

\section*{Risks}
Venipuncture only.

\section*{What Happens After the Test}
Low levels prompt investigation and correction of fat malabsorption; high-dose supplementation prescribed in deficiency. Treatment reverses neurologic deficits in early stages; advanced neurologic injury may be irreversible.

\section*{Understanding Results}

\subsection*{Below Normal}
\begin{itemize}
  \item Fat malabsorption: cystic fibrosis, cholestatic liver disease (especially children with chronic cholestasis), short bowel
  \item Abetalipoproteinemia: hereditary inability to assemble apo-B-containing lipoproteins
  \item Premature infants: born with low stores
  \item Long-term inadequate parenteral nutrition
  \item Severe protein-energy malnutrition
  \item Malabsorption syndromes (Crohn's, celiac disease)
  \item Symptoms: ataxia, peripheral neuropathy, retinopathy, myopathy, hemolytic anemia in neonates
\end{itemize}

\subsection*{Above Normal}
\begin{itemize}
  \item Excess supplementation: high-dose vitamin E ($>$ 1,000 mg/day) increases risk of bleeding (interference with vitamin K); meta-analyses suggested all-cause mortality risk at high doses
  \item Parenteral nutrition with excessive supplementation
  \item Isolated high intake without toxicity is rare due to urinary excretion
\end{itemize}

\section*{Normal Results}
\begin{itemize}
  \item \textbf{Serum alpha-tocopherol (adult):} 5.5--17 mg/L (12--42 $\mu$mol/L)
  \item \textbf{Children:} 3--18 mg/L (7--42 $\mu$mol/L)
  \item \textbf{Alpha-tocopherol/total lipid ratio:} $>$ 0.8 mg/g lipid (consistent adequacy)
  \item \textbf{Deficiency:} $<$ 5 mg/L; $<$ 0.8 mg/g corrected is biochemical
  \item \textbf{Recommended daily intake:} Adults 15 mg/day; pregnant/lactating women 19 mg/day
\end{itemize}

\section*{Follow-up Tests}
\begin{itemize}
  \item \textbf{Serum lipids (total cholesterol, triglycerides):} Allow correction for lipid-standardized ratio
  \item \textbf{Other fat-soluble vitamins (A, D, K)}: Co-existent deficiencies common in fat malabsorption
  \item \textbf{72-hour fecal fat and stool elastase}: Diagnose fat malabsorption/pancreatic insufficiency
  \item \textbf{Prothrombin time (PT/INR)}: Vitamin K deficiency due to fat malabsorption
  \item \textbf{Complete blood count and peripheral smear}: Hemolytic anemia in neonates
  \item \textbf{Neurological examination and MRI}: Spinocerebellar atrophy with chronic deficiency
  \item \textbf{Genetic testing (MTTP gene)}: If abetalipoproteinemia suspected
\end{itemize}

\section*{References}
\begin{enumerate}
  \item Traber MG, Stevens JF. Vitamin E: mechanisms of antioxidant action. Annu Rev Nutr. 2011;31:159-178.
  \item Institute of Medicine. Dietary Reference Intakes for Vitamin C, Vitamin E, Selenium, and Carotenoids. National Academies Press; 2000.
  \item Sokol RJ. Vitamin E deficiency and neurologic disease. Annu Rev Nutr. 1988;8:349-367.
  \item Schmölz L, et al. Vitamin E: an updated review. Mol Nutr Food Res. 2016;60(1):45-63.
\end{enumerate}

\clearpage
\section*{Regulatory Status and Authorization Date}
Regulatory status applies to the specific assay, instrument, manufacturer, and intended use rather than to the test name alone. No single approval date can be assigned to this test category. For a commercial assay, verify the current product in the relevant regulator's database: the U.S. Food and Drug Administration (FDA) clearance, approval, or authorization; India's Central Drugs Standard Control Organisation (CDSCO) licence or permission; the European Union In Vitro Diagnostic Regulation (EU IVDR) CE certificate issued through the applicable conformity-assessment route; or the United Kingdom Medicines and Healthcare products Regulatory Agency (MHRA) registration and UKCA/CE status. Laboratory-developed tests may not have an individual FDA, CDSCO, EU IVDR, or MHRA approval date.
\clearpage
